% from main.tex
%%%%%%%%%%  scan idx 113  |  folio 100  %%%%%%%%%%
vraie » flèche de \(\Cpx^-(\mathcal C_S)\). L'assertion résulte alors des définitions et de la formule \(C(u)^i=E^{i+1}\oplus F^i\), où \(C(u)\) désigne le cône de \(u\).

\textbf{Scholie 2.6.} Par « rotation » des triangles, on voit que des propriétés de pseudo-cohérence sur deux des objets du triangle \((*)\) en impliquent une autre sur le troisième. On peut résumer les différents cas dans le tableau ci-dessous (les petites lettres désignent des degrés de pseudo-cohérence et chaque ligne correspond à une implication dont le but est souligné) :
\[
\begin{array}{ccc}
E&F&G\\[0.2em]
n+1&n&\underline n\\
n&\underline n&n\\
\underline n&n&n-1 .
\end{array}
\]

Nous noterons \(\Dpx(\mathcal C)_{\mathcal C_0\text{-}n\text{-coh}}\) (resp. \(\Dpx(\mathcal C)_{\mathcal C_0\text{-}\mathrm{coh}}\)) la sous-catégorie fibrée strictement pleine de \(\Dpx(\mathcal C)\) dont la fibre en chaque \(X\in\Ob S\) est formée des objets de \(\Dpx(\mathcal C_X)\) qui sont \(n\)-pseudo-cohérents (resp. pseudo-cohérents) relativement à \(\mathcal C_0\). La catégorie \(\Dpx(\mathcal C)_{\mathcal C_0\text{-}\mathrm{coh}}\) est, en vertu de (2.5), une \emph{sous-\(S\)-catégorie triangulée} de \(\Dpx(\mathcal C)\) (1.1).

\textbf{Proposition 2.7.} Supposons que \(S\) soit un site ponctuel, donc \(\mathcal C=\mathcal C_S\), \(\mathcal C_0=\mathcal C_{0,S}\). Alors :

\textbf{a)} Si \(F\in\Ob\Dpx(\mathcal C)\) est pseudo-cohérent, il existe un quasi-isomorphisme \(E\longrightarrow F\), avec \(E\) strictement pseudo-cohérent (2.1).

\textbf{b)} Désignons par \(\Kpx^{-,b}(\mathcal C_0)\) la sous-catégorie triangulée pleine de \(\Kpx^-(\mathcal C_0)\) formée des complexes acycliques. Le foncteur canonique
\[
  \Kpx^-(\mathcal C_0)/\Kpx^{-,b}(\mathcal C_0)
  \longrightarrow \Dpx(\mathcal C)
\]
